\IfFileExists{stacks-project.cls}{%
\documentclass{stacks-project}
}{%
\documentclass{amsart}
}

% For dealing with references we use the comment environment
\usepackage{verbatim}
\newenvironment{reference}{\comment}{\endcomment}
%\newenvironment{reference}{}{}
\newenvironment{slogan}{\comment}{\endcomment}
\newenvironment{history}{\comment}{\endcomment}

% For commutative diagrams we use Xy-pic
\usepackage[all]{xy}

% We use 2cell for 2-commutative diagrams.
\xyoption{2cell}
\UseAllTwocells

% We use multicol for the list of chapters between chapters
\usepackage{multicol}

% This is generally recommended for better output
\usepackage{lmodern}
\usepackage[T1]{fontenc}

% For cross-file-references
\usepackage{xr-hyper}

% Package for hypertext links:
\usepackage{hyperref}

% For any local file, say "hello.tex" you want to link to please
% use \externaldocument[hello-]{hello}
\externaldocument[introduction-]{introduction}
\externaldocument[conventions-]{conventions}
\externaldocument[sets-]{sets}
\externaldocument[categories-]{categories}
\externaldocument[topology-]{topology}
\externaldocument[sheaves-]{sheaves}
\externaldocument[sites-]{sites}
\externaldocument[stacks-]{stacks}
\externaldocument[fields-]{fields}
\externaldocument[algebra-]{algebra}
\externaldocument[brauer-]{brauer}
\externaldocument[homology-]{homology}
\externaldocument[derived-]{derived}
\externaldocument[simplicial-]{simplicial}
\externaldocument[more-algebra-]{more-algebra}
\externaldocument[smoothing-]{smoothing}
\externaldocument[modules-]{modules}
\externaldocument[sites-modules-]{sites-modules}
\externaldocument[injectives-]{injectives}
\externaldocument[cohomology-]{cohomology}
\externaldocument[sites-cohomology-]{sites-cohomology}
\externaldocument[dga-]{dga}
\externaldocument[dpa-]{dpa}
\externaldocument[sdga-]{sdga}
\externaldocument[hypercovering-]{hypercovering}
\externaldocument[schemes-]{schemes}
\externaldocument[constructions-]{constructions}
\externaldocument[properties-]{properties}
\externaldocument[morphisms-]{morphisms}
\externaldocument[coherent-]{coherent}
\externaldocument[divisors-]{divisors}
\externaldocument[limits-]{limits}
\externaldocument[varieties-]{varieties}
\externaldocument[topologies-]{topologies}
\externaldocument[descent-]{descent}
\externaldocument[perfect-]{perfect}
\externaldocument[more-morphisms-]{more-morphisms}
\externaldocument[flat-]{flat}
\externaldocument[groupoids-]{groupoids}
\externaldocument[more-groupoids-]{more-groupoids}
\externaldocument[etale-]{etale}
\externaldocument[chow-]{chow}
\externaldocument[intersection-]{intersection}
\externaldocument[pic-]{pic}
\externaldocument[weil-]{weil}
\externaldocument[adequate-]{adequate}
\externaldocument[dualizing-]{dualizing}
\externaldocument[duality-]{duality}
\externaldocument[discriminant-]{discriminant}
\externaldocument[derham-]{derham}
\externaldocument[local-cohomology-]{local-cohomology}
\externaldocument[algebraization-]{algebraization}
\externaldocument[curves-]{curves}
\externaldocument[resolve-]{resolve}
\externaldocument[models-]{models}
\externaldocument[functors-]{functors}
\externaldocument[equiv-]{equiv}
\externaldocument[pione-]{pione}
\externaldocument[etale-cohomology-]{etale-cohomology}
\externaldocument[proetale-]{proetale}
\externaldocument[relative-cycles-]{relative-cycles}
\externaldocument[more-etale-]{more-etale}
\externaldocument[trace-]{trace}
\externaldocument[crystalline-]{crystalline}
\externaldocument[spaces-]{spaces}
\externaldocument[spaces-properties-]{spaces-properties}
\externaldocument[spaces-morphisms-]{spaces-morphisms}
\externaldocument[decent-spaces-]{decent-spaces}
\externaldocument[spaces-cohomology-]{spaces-cohomology}
\externaldocument[spaces-limits-]{spaces-limits}
\externaldocument[spaces-divisors-]{spaces-divisors}
\externaldocument[spaces-over-fields-]{spaces-over-fields}
\externaldocument[spaces-topologies-]{spaces-topologies}
\externaldocument[spaces-descent-]{spaces-descent}
\externaldocument[spaces-perfect-]{spaces-perfect}
\externaldocument[spaces-more-morphisms-]{spaces-more-morphisms}
\externaldocument[spaces-flat-]{spaces-flat}
\externaldocument[spaces-groupoids-]{spaces-groupoids}
\externaldocument[spaces-more-groupoids-]{spaces-more-groupoids}
\externaldocument[bootstrap-]{bootstrap}
\externaldocument[spaces-pushouts-]{spaces-pushouts}
\externaldocument[spaces-chow-]{spaces-chow}
\externaldocument[groupoids-quotients-]{groupoids-quotients}
\externaldocument[spaces-more-cohomology-]{spaces-more-cohomology}
\externaldocument[spaces-simplicial-]{spaces-simplicial}
\externaldocument[spaces-duality-]{spaces-duality}
\externaldocument[formal-spaces-]{formal-spaces}
\externaldocument[restricted-]{restricted}
\externaldocument[spaces-resolve-]{spaces-resolve}
\externaldocument[formal-defos-]{formal-defos}
\externaldocument[defos-]{defos}
\externaldocument[cotangent-]{cotangent}
\externaldocument[examples-defos-]{examples-defos}
\externaldocument[algebraic-]{algebraic}
\externaldocument[examples-stacks-]{examples-stacks}
\externaldocument[stacks-sheaves-]{stacks-sheaves}
\externaldocument[criteria-]{criteria}
\externaldocument[artin-]{artin}
\externaldocument[quot-]{quot}
\externaldocument[stacks-properties-]{stacks-properties}
\externaldocument[stacks-morphisms-]{stacks-morphisms}
\externaldocument[stacks-limits-]{stacks-limits}
\externaldocument[stacks-cohomology-]{stacks-cohomology}
\externaldocument[stacks-perfect-]{stacks-perfect}
\externaldocument[stacks-introduction-]{stacks-introduction}
\externaldocument[stacks-more-morphisms-]{stacks-more-morphisms}
\externaldocument[stacks-geometry-]{stacks-geometry}
\externaldocument[moduli-]{moduli}
\externaldocument[moduli-curves-]{moduli-curves}
\externaldocument[examples-]{examples}
\externaldocument[exercises-]{exercises}
\externaldocument[guide-]{guide}
\externaldocument[desirables-]{desirables}
\externaldocument[coding-]{coding}
\externaldocument[obsolete-]{obsolete}
\externaldocument[fdl-]{fdl}
\externaldocument[index-]{index}

% Theorem environments.
%
\theoremstyle{plain}
\newtheorem{theorem}[subsection]{Theorem}
\newtheorem{proposition}[subsection]{Proposition}
\newtheorem{lemma}[subsection]{Lemma}

\theoremstyle{definition}
\newtheorem{definition}[subsection]{Definition}
\newtheorem{example}[subsection]{Example}
\newtheorem{exercise}[subsection]{Exercise}
\newtheorem{situation}[subsection]{Situation}

\theoremstyle{remark}
\newtheorem{remark}[subsection]{Remark}
\newtheorem{remarks}[subsection]{Remarks}

\numberwithin{equation}{subsection}

% Macros
%
\def\lim{\mathop{\mathrm{lim}}\nolimits}
\def\colim{\mathop{\mathrm{colim}}\nolimits}
\def\Spec{\mathop{\mathrm{Spec}}}
\def\Hom{\mathop{\mathrm{Hom}}\nolimits}
\def\Ext{\mathop{\mathrm{Ext}}\nolimits}
\def\SheafHom{\mathop{\mathcal{H}\!\mathit{om}}\nolimits}
\def\SheafExt{\mathop{\mathcal{E}\!\mathit{xt}}\nolimits}
\def\Sch{\mathit{Sch}}
\def\Mor{\mathop{\mathrm{Mor}}\nolimits}
\def\Ob{\mathop{\mathrm{Ob}}\nolimits}
\def\Sh{\mathop{\mathit{Sh}}\nolimits}
\def\NL{\mathop{N\!L}\nolimits}
\def\CH{\mathop{\mathrm{CH}}\nolimits}
\def\proetale{{pro\text{-}\acute{e}tale}}
\def\etale{{\acute{e}tale}}
\def\QCoh{\mathit{QCoh}}
\def\Ker{\mathop{\mathrm{Ker}}}
\def\Im{\mathop{\mathrm{Im}}}
\def\Coker{\mathop{\mathrm{Coker}}}
\def\Coim{\mathop{\mathrm{Coim}}}

% Boxtimes
%
\DeclareMathSymbol{\boxtimes}{\mathbin}{AMSa}{"02}

%
% Macros for moduli stacks/spaces
%
\def\QCohstack{\mathcal{QC}\!\mathit{oh}}
\def\Cohstack{\mathcal{C}\!\mathit{oh}}
\def\Spacesstack{\mathcal{S}\!\mathit{paces}}
\def\Quotfunctor{\mathrm{Quot}}
\def\Hilbfunctor{\mathrm{Hilb}}
\def\Curvesstack{\mathcal{C}\!\mathit{urves}}
\def\Polarizedstack{\mathcal{P}\!\mathit{olarized}}
\def\Complexesstack{\mathcal{C}\!\mathit{omplexes}}
% \Pic is the operator that assigns to X its picard group, usage \Pic(X)
% \Picardstack_{X/B} denotes the Picard stack of X over B
% \Picardfunctor_{X/B} denotes the Picard functor of X over B
\def\Pic{\mathop{\mathrm{Pic}}\nolimits}
\def\Picardstack{\mathcal{P}\!\mathit{ic}}
\def\Picardfunctor{\mathrm{Pic}}
\def\Deformationcategory{\mathcal{D}\!\mathit{ef}}
\usepackage{cleveref}
\crefname{originalreference}{original reference}{original references}
\Crefname{originalreference}{Original reference}{Original references}
\makeatletter
\newlabel{fields-lemma-normal-closure}{{09DT}{source}{}{sourcepointer-1}{}}
\newlabel{fields-lemma-normal-closure@cref}{{[originalreference][0][]09DT}{[1][0][]source}{}{}{}}
\makeatother
\crefname{originalreference}{original reference}{original references}
\Crefname{originalreference}{Original reference}{Original references}
\makeatletter
\newlabel{groupoids-lemma-compact-set-in-affine}{{0B7U}{source}{}{sourcepointer-2}{}}
\newlabel{groupoids-lemma-compact-set-in-affine@cref}{{[originalreference][0][]0B7U}{[1][0][]source}{}{}{}}
\makeatother
\crefname{originallemma}{original lemma}{original lemmas}
\Crefname{originallemma}{Original lemma}{Original lemmas}
\makeatletter
\newlabel{lemma-group-space-scheme-locally-finite-type-over-k}{{0B8F}{source}{}{sourcepointer-3}{}}
\newlabel{lemma-group-space-scheme-locally-finite-type-over-k@cref}{{[originallemma][0][]0B8F}{[1][0][]source}{}{}{}}
\makeatother
\crefname{originallemma}{original lemma}{original lemmas}
\Crefname{originallemma}{Original lemma}{Original lemmas}
\makeatletter
\newlabel{lemma-group-space-scheme-over-kbar}{{0B8E}{source}{}{sourcepointer-4}{}}
\newlabel{lemma-group-space-scheme-over-kbar@cref}{{[originallemma][0][]0B8E}{[1][0][]source}{}{}{}}
\makeatother
\crefname{originalreference}{original reference}{original references}
\Crefname{originalreference}{Original reference}{Original references}
\makeatletter
\newlabel{limits-lemma-descend-opens}{{01Z4}{source}{}{sourcepointer-5}{}}
\newlabel{limits-lemma-descend-opens@cref}{{[originalreference][0][]01Z4}{[1][0][]source}{}{}{}}
\makeatother
\crefname{originalreference}{original reference}{original references}
\Crefname{originalreference}{Original reference}{Original references}
\makeatletter
\newlabel{limits-lemma-limit-affine}{{01Z6}{source}{}{sourcepointer-6}{}}
\newlabel{limits-lemma-limit-affine@cref}{{[originalreference][0][]01Z6}{[1][0][]source}{}{}{}}
\makeatother
\crefname{originalreference}{original reference}{original references}
\Crefname{originalreference}{Original reference}{Original references}
\makeatletter
\newlabel{spaces-over-fields-lemma-base-change-by-Galois}{{0B86}{source}{}{sourcepointer-7}{}}
\newlabel{spaces-over-fields-lemma-base-change-by-Galois@cref}{{[originalreference][0][]0B86}{[1][0][]source}{}{}{}}
\makeatother
\crefname{originalreference}{original reference}{original references}
\Crefname{originalreference}{Original reference}{Original references}
\makeatletter
\newlabel{spaces-over-fields-lemma-scheme-after-purely-inseparable-base-change}{{0B83}{source}{}{sourcepointer-8}{}}
\newlabel{spaces-over-fields-lemma-scheme-after-purely-inseparable-base-change@cref}{{[originalreference][0][]0B83}{[1][0][]source}{}{}{}}
\makeatother
\crefname{originalreference}{original reference}{original references}
\Crefname{originalreference}{Original reference}{Original references}
\makeatletter
\newlabel{spaces-over-fields-lemma-scheme-after-purely-transcendental-base-change}{{0B85}{source}{}{sourcepointer-9}{}}
\newlabel{spaces-over-fields-lemma-scheme-after-purely-transcendental-base-change@cref}{{[originalreference][0][]0B85}{[1][0][]source}{}{}{}}
\makeatother
\crefname{originalreference}{original reference}{original references}
\Crefname{originalreference}{Original reference}{Original references}
\makeatletter
\newlabel{spaces-over-fields-lemma-scheme-over-algebraic-closure-enough-affines}{{0B88}{source}{}{sourcepointer-10}{}}
\newlabel{spaces-over-fields-lemma-scheme-over-algebraic-closure-enough-affines@cref}{{[originalreference][0][]0B88}{[1][0][]source}{}{}{}}
\makeatother
\crefname{originalreference}{original reference}{original references}
\Crefname{originalreference}{Original reference}{Original references}
\makeatletter
\newlabel{spaces-over-fields-lemma-when-quotient-scheme-at-point}{{0B87}{source}{}{sourcepointer-11}{}}
\newlabel{spaces-over-fields-lemma-when-quotient-scheme-at-point@cref}{{[originalreference][0][]0B87}{[1][0][]source}{}{}{}}
\makeatother
\crefname{originalreference}{original reference}{original references}
\Crefname{originalreference}{Original reference}{Original references}
\makeatletter
\newlabel{spaces-over-fields-lemma-when-scheme-after-base-change}{{0B84}{source}{}{sourcepointer-12}{}}
\newlabel{spaces-over-fields-lemma-when-scheme-after-base-change@cref}{{[originalreference][0][]0B84}{[1][0][]source}{}{}{}}
\makeatother
\crefname{originalreference}{original reference}{original references}
\Crefname{originalreference}{Original reference}{Original references}
\makeatletter
\newlabel{varieties-lemma-make-Jacobson}{{0479}{source}{}{sourcepointer-13}{}}
\newlabel{varieties-lemma-make-Jacobson@cref}{{[originalreference][0][]0479}{[1][0][]source}{}{}{}}
\makeatother
\title{Separated group spaces over a field are schemes}
\author{The Stacks Project Authors}
\date{stacks-project:spaces-more-groupoids.tex:local-snapshot-20261012; GNU FDL 1.2 or later}
\begin{document}
\maketitle
\noindent\textit{Editorial note.} The author flags the sibling proof (spaces-over-fields lemma used for the lft case) as recommended reading; the proof itself is complete in-span with no omission markers. Cross-chapter same-name note: no same-label card exists in any harvested chapter; the proposition is the general form of the locally-finite-type case carded at L26664's sibling batch position (lemma-group-space-scheme-locally-finite-type-over-k was rejected as a pure three-way reference list, so no double count arises).\par\medskip
\noindent\textit{Allowed standard prerequisites (editorial).}\begin{itemize}
\item Group space over the algebraic closure. lemma-group-space-scheme-over-kbar (same chapter, carded L26646). G\_\{k-bar\} is a scheme, the anchor for the descent chain.
\item When a space becomes a scheme after base change. Spaces over Fields lemmas scheme-after-purely-transcendental-base-change, scheme-after-purely-inseparable-base-change, when-scheme-after-base-change, base-change-by-Galois, when-quotient-scheme-at-point (cross-chapter, named block). The four-step field-extension chain and the Galois descent that closes the proof.
\item Limits of affine opens. Limits, lemmas descend-opens and limit-affine (cross-chapter, named). Recover W as the base change of an affine open U' over a finite subextension.
\end{itemize}\medskip
\ifdefined\hypertarget\hypertarget{sourcepointer-1}{}\fi
\noindent\textit{Source reference.} Original reference 09DT: focused-seeds/geometry/S9-smg-group-space-scheme-over-field/ref-tags.txt, source lines 5--5. Navigation pointer only: the counted proof cites this statement by \textbackslash{}ref as a standing input; the deduction the counted proof has to make is carried out inside the counted span, and no part of that deduction is delegated to this pointer.\par\medskip
\ifdefined\hypertarget\hypertarget{sourcepointer-2}{}\fi
\noindent\textit{Source reference.} Original reference 0B7U: focused-seeds/geometry/S9-smg-group-space-scheme-over-field/ref-tags.txt, source lines 6--6. Navigation pointer only: the counted proof cites this statement by \textbackslash{}ref as a standing input; the deduction the counted proof has to make is carried out inside the counted span, and no part of that deduction is delegated to this pointer.\par\medskip
\ifdefined\hypertarget\hypertarget{sourcepointer-3}{}\fi
\noindent\textit{Source reference.} Original lemma 0B8F: sources/stacks-spaces-more-groupoids/pinned-original-source.tex, source lines 1296--1296. Navigation pointer only: the counted proof cites this statement by \textbackslash{}ref as a standing input; the deduction the counted proof has to make is carried out inside the counted span, and no part of that deduction is delegated to this pointer.\par\medskip
\ifdefined\hypertarget\hypertarget{sourcepointer-4}{}\fi
\noindent\textit{Source reference.} Original lemma 0B8E: sources/stacks-spaces-more-groupoids/pinned-original-source.tex, source lines 1258--1258. Navigation pointer only: the counted proof cites this statement by \textbackslash{}ref as a standing input; the deduction the counted proof has to make is carried out inside the counted span, and no part of that deduction is delegated to this pointer.\par\medskip
\ifdefined\hypertarget\hypertarget{sourcepointer-5}{}\fi
\noindent\textit{Source reference.} Original reference 01Z4: focused-seeds/geometry/S9-smg-group-space-scheme-over-field/ref-tags.txt, source lines 7--7. Navigation pointer only: the counted proof cites this statement by \textbackslash{}ref as a standing input; the deduction the counted proof has to make is carried out inside the counted span, and no part of that deduction is delegated to this pointer.\par\medskip
\ifdefined\hypertarget\hypertarget{sourcepointer-6}{}\fi
\noindent\textit{Source reference.} Original reference 01Z6: focused-seeds/geometry/S9-smg-group-space-scheme-over-field/ref-tags.txt, source lines 8--8. Navigation pointer only: the counted proof cites this statement by \textbackslash{}ref as a standing input; the deduction the counted proof has to make is carried out inside the counted span, and no part of that deduction is delegated to this pointer.\par\medskip
\ifdefined\hypertarget\hypertarget{sourcepointer-7}{}\fi
\noindent\textit{Source reference.} Original reference 0B86: focused-seeds/geometry/S9-smg-group-space-scheme-over-field/ref-tags.txt, source lines 9--9. Navigation pointer only: the counted proof cites this statement by \textbackslash{}ref as a standing input; the deduction the counted proof has to make is carried out inside the counted span, and no part of that deduction is delegated to this pointer.\par\medskip
\ifdefined\hypertarget\hypertarget{sourcepointer-8}{}\fi
\noindent\textit{Source reference.} Original reference 0B83: focused-seeds/geometry/S9-smg-group-space-scheme-over-field/ref-tags.txt, source lines 10--10. Navigation pointer only: the counted proof cites this statement by \textbackslash{}ref as a standing input; the deduction the counted proof has to make is carried out inside the counted span, and no part of that deduction is delegated to this pointer.\par\medskip
\ifdefined\hypertarget\hypertarget{sourcepointer-9}{}\fi
\noindent\textit{Source reference.} Original reference 0B85: focused-seeds/geometry/S9-smg-group-space-scheme-over-field/ref-tags.txt, source lines 11--11. Navigation pointer only: the counted proof cites this statement by \textbackslash{}ref as a standing input; the deduction the counted proof has to make is carried out inside the counted span, and no part of that deduction is delegated to this pointer.\par\medskip
\ifdefined\hypertarget\hypertarget{sourcepointer-10}{}\fi
\noindent\textit{Source reference.} Original reference 0B88: focused-seeds/geometry/S9-smg-group-space-scheme-over-field/ref-tags.txt, source lines 12--12. Navigation pointer only: the counted proof cites this statement by \textbackslash{}ref as a standing input; the deduction the counted proof has to make is carried out inside the counted span, and no part of that deduction is delegated to this pointer.\par\medskip
\ifdefined\hypertarget\hypertarget{sourcepointer-11}{}\fi
\noindent\textit{Source reference.} Original reference 0B87: focused-seeds/geometry/S9-smg-group-space-scheme-over-field/ref-tags.txt, source lines 13--13. Navigation pointer only: the counted proof cites this statement by \textbackslash{}ref as a standing input; the deduction the counted proof has to make is carried out inside the counted span, and no part of that deduction is delegated to this pointer.\par\medskip
\ifdefined\hypertarget\hypertarget{sourcepointer-12}{}\fi
\noindent\textit{Source reference.} Original reference 0B84: focused-seeds/geometry/S9-smg-group-space-scheme-over-field/ref-tags.txt, source lines 14--14. Navigation pointer only: the counted proof cites this statement by \textbackslash{}ref as a standing input; the deduction the counted proof has to make is carried out inside the counted span, and no part of that deduction is delegated to this pointer.\par\medskip
\ifdefined\hypertarget\hypertarget{sourcepointer-13}{}\fi
\noindent\textit{Source reference.} Original reference 0479: focused-seeds/geometry/S9-smg-group-space-scheme-over-field/ref-tags.txt, source lines 15--15. Navigation pointer only: the counted proof cites this statement by \textbackslash{}ref as a standing input; the deduction the counted proof has to make is carried out inside the counted span, and no part of that deduction is delegated to this pointer.\par\medskip
\begin{proposition}
\label{proposition-group-space-scheme-over-field}
Let $k$ be a field. Let $G$ be a group algebraic space over $k$.
If $G$ is separated, then $G$ is a scheme.
\end{proposition}

\begin{proof}
This lemma generalizes
Lemma \ref{lemma-group-space-scheme-locally-finite-type-over-k}
(which covers all cases one cares about in practice).
The proof is very similar to the proof of
Spaces over Fields, Lemma
\ref{spaces-over-fields-lemma-scheme-over-algebraic-closure-enough-affines}
used in the proof of
Lemma \ref{lemma-group-space-scheme-locally-finite-type-over-k}
and we encourage the reader to read that proof first.

\medskip\noindent
By Lemma \ref{lemma-group-space-scheme-over-kbar} the base
change $G_{\overline{k}}$ is a scheme.
Let $K/k$ be a purely transcendental extension of very large
transcendence degree. By Spaces over Fields, Lemma
\ref{spaces-over-fields-lemma-scheme-after-purely-transcendental-base-change}
it suffices to show that $G_K$ is a scheme.
Let $K^{perf}$ be the perfect closure of $K$. By
Spaces over Fields, Lemma
\ref{spaces-over-fields-lemma-scheme-after-purely-inseparable-base-change}
it suffices to show that $G_{K^{perf}}$ is a scheme.
Let $K \subset K^{perf} \subset \overline{K}$ be the algebraic closure
of $K$. We may choose an embedding $\overline{k} \to \overline{K}$ over
$k$, so that $G_{\overline{K}}$ is the base change of the scheme
$G_{\overline{k}}$ by $\overline{k} \to \overline{K}$. By
Varieties, Lemma \ref{varieties-lemma-make-Jacobson}
we see that $G_{\overline{K}}$ is a Jacobson scheme all of whose
closed points have residue field $\overline{K}$.

\medskip\noindent
Since $G_{\overline{K}} \to G_{K^{perf}}$ is surjective, it suffices to
show that the image $g \in |G_{K^{perf}}|$ of an arbitrary closed point of
$G_{\overline{K}}$ is in the schematic locus of $G_K$.
In particular, we may represent $g$ by a morphism
$g : \Spec(L) \to G_{K^{perf}}$ where $L/K^{perf}$ is separable algebraic
(for example we can take $L = \overline{K}$). Thus the scheme
\begin{align*}
T & = \Spec(L) \times_{G_{K^{perf}}} G_{\overline{K}} \\
& =
\Spec(L) \times_{\Spec(K^{perf})} \Spec(\overline{K}) \\
& =
\Spec(L \otimes_{K^{perf}} \overline{K})
\end{align*}
is the spectrum of a $\overline{K}$-algebra which is a filtered colimit
of algebras which are finite products of copies of $\overline{K}$.
Thus by Groupoids, Lemma \ref{groupoids-lemma-compact-set-in-affine}
we can find an affine open $W \subset G_{\overline{K}}$ containing
the image of $g_{\overline{K}} : T \to G_{\overline{K}}$.

\medskip\noindent
Choose a quasi-compact open $V \subset G_{K^{perf}}$ containing the
image of $W$. By Spaces over Fields,
Lemma \ref{spaces-over-fields-lemma-when-scheme-after-base-change}
we see that $V_{K'}$ is a scheme for some finite extension $K'/K^{perf}$.
After enlarging $K'$ we may assume that there exists an affine open
$U' \subset V_{K'} \subset G_{K'}$ whose base change to $\overline{K}$
recovers $W$ (use that $V_{\overline{K}}$ is the limit of the schemes
$V_{K''}$ for $K' \subset K'' \subset \overline{K}$ finite and use
Limits, Lemmas \ref{limits-lemma-descend-opens} and
\ref{limits-lemma-limit-affine}). We may assume
that $K'/K^{perf}$ is a Galois extension (take the normal closure
Fields, Lemma \ref{fields-lemma-normal-closure} and use
that $K^{perf}$ is perfect). Set $H = \text{Gal}(K'/K^{perf})$.
By construction the $H$-invariant closed subscheme
$\Spec(L) \times_{G_{K^{perf}}} G_{K'}$ is contained in $U'$.
By
Spaces over Fields, Lemmas
\ref{spaces-over-fields-lemma-base-change-by-Galois} and
\ref{spaces-over-fields-lemma-when-quotient-scheme-at-point} we conclude.
\end{proof}

\end{document}
